\documentclass[12pt]{article}
\usepackage{amsmath,amsfonts,amssymb}
\usepackage{geometry}
\geometry{a4paper,margin=1in}

%%
%%  This document presents a detailed derivation of the minimum–impulse
%%  transfer between two nearby low-eccentricity orbits in three
%%  dimensions.  The derivation follows the approach of Edelbaum and
%%  collaborators but fills in intermediate steps and corrects several
%%  typographical errors that appear in the scanned article.  Throughout
%%  the derivation we adopt Lagrange's planetary variables and linearize
%%  the equations of motion about a circular reference orbit of
%%  semi-major axis a_0.
%%

\title{Detailed derivation of minimum-impulse transfers near a circular orbit}
\author{}
\date{}

\begin{document}
\maketitle

\section*{1\quad Introduction}

Optimal guidance of spacecraft frequently reduces to finding a set of
impulsive velocity changes that transforms one orbit into another
while minimising the total velocity expenditure.  When the orbits of
interest are close to a common circular orbit, the problem may be
linearised with respect to small deviations of the orbital elements.
The resulting linear optimal control problem is amenable to analytic
solution and yields insight into the number, timing and direction of
optimal impulses.  In what follows we derive the complete solution of
the minimum-impulse problem for transfers between arbitrary nearby
orbits with small eccentricities and small mutual inclination.

\paragraph{State variables.}  The departure and arrival orbits are
\emph{described} by deviations in the semi-major axis $a$, the eccentricity
vector $(e_x,e_y)$ and the inclination vector $(i_x,i_y)$ relative
to the reference circular orbit of semi-major axis~$a_0$.  For small
eccentricity $e\ll 1$ and small inclination $i\ll 1$ the Cartesian
components $e_x$, $e_y$, $i_x$ and $i_y$ are proportional to the usual
scalar eccentricity $e$ and inclination $i$, respectively.  The
state vector
\[
  x = \begin{pmatrix}
    x_1 \\ x_2 \\ x_3 \\ x_4 \\ x_5
  \end{pmatrix}
  = \begin{pmatrix}
    \displaystyle \frac{a-a_0}{a_0} \\ e_y \\ e_x \\ i_y \\ i_x
  \end{pmatrix}
\]
will be used in the sequel, together with the independent variable
\emph{fuel} $u$, defined as the time integral of the thrust
acceleration.  Because the thrust is assumed to be arbitrarily large
(impulsive), $u$ is simply the sum of the magnitudes of the velocity
impulses.

\paragraph{Control variables.}  An impulse of magnitude~$\Delta v$ is
characterised by the direction cosines $l_R$, $l_T$ and~$l_N$ of the
thrust vector along the radial, circumferential and normal directions
of the reference orbit, respectively.  For an infinitesimal change of
magnitude $\mathrm{d}u$ the corresponding change in the spacecraft
velocity is $\mathrm{d}\mathbf{v} = (l_R, l_T, l_N)\,\mathrm{d}u$.

\paragraph{Planar equations of motion.}  To first order in the
perturbations the equations of motion follow from Lagrange's planetary
equations.  Linearisation about the reference circular orbit yields
\begin{align}
  \frac{\mathrm{d}x_1}{\mathrm{d}u} &= 2\,\sqrt{\frac{a_0}{K}}\,l_T,
    \tag{1}\label{eq:dx1du}\\
  \frac{\mathrm{d}x_2}{\mathrm{d}u} &= \sqrt{\frac{a_0}{K}}
    \bigl(2\,l_T\sin\theta - l_R\cos\theta\bigr),
    \tag{2}\label{eq:dx2du}\\
  \frac{\mathrm{d}x_3}{\mathrm{d}u} &= \sqrt{\frac{a_0}{K}}
    \bigl(2\,l_T\cos\theta + l_R\sin\theta\bigr),
    \tag{3}\label{eq:dx3du}\\
  \frac{\mathrm{d}x_4}{\mathrm{d}u} &= \sqrt{\frac{a_0}{K}}\,l_N\sin\theta,
    \tag{4}\label{eq:dx4du}\\
  \frac{\mathrm{d}x_5}{\mathrm{d}u} &= \sqrt{\frac{a_0}{K}}\,l_N\cos\theta.
    \tag{5}\label{eq:dx5du}
\end{align}
Here $K$ denotes the gravitational constant $GM$ of the central body
and $\theta$ is the true longitude (central angle) of the spacecraft
with respect to a fixed reference direction in the reference orbital
plane.  Equations~\eqref{eq:dx1du}--\eqref{eq:dx5du} are easily
derived from the Gauss variational equations: a purely tangential
impulse changes the semi-major axis through an energy change,
whereas radial and tangential impulses alter the components of the
eccentricity vector, and a normal impulse changes the inclination.

\section*{2\quad Hamiltonian formulation and primer vector}

We convert the optimal impulsive control problem into a static
Hamiltonian system by applying Pontryagin's maximum principle.  The
Hamiltonian is linear in the control because the state equations are
linear in the thrust direction cosines,
\begin{equation}
  H = \sum_{i=1}^5 \lambda_i\,\frac{\mathrm{d}x_i}{\mathrm{d}u},
  \tag{6}\label{eq:Hamiltonian}
\end{equation}
where $\lambda_i$ are the adjoint variables (Lagrange multipliers) and
the explicit factor $p = \max\|\lambda\|$ is the magnitude of the
\emph{primer vector} introduced by Lawden.  Writing the components of
the primer vector in the radial, tangential and normal directions as
$\lambda$, $\mu$ and $\nu$, respectively, we have
\begin{equation}
  l_R = \frac{\lambda}{p},\quad l_T = \frac{\mu}{p},\quad
  l_N = \frac{\nu}{p},
  \tag{7}\label{eq:thrustdircos}
\end{equation}
so that an optimal impulse points in the direction of the primer
vector.  Substituting the state equations
\eqref{eq:dx1du}--\eqref{eq:dx5du} into the Hamiltonian yields the
components of the primer vector in terms of the adjoint variables
$\lambda_i$:
\begin{align}
  \lambda &= \sqrt{\frac{a_0}{K}}
    \bigl(-\lambda_2\cos\theta + \lambda_3\sin\theta\bigr),
    \tag{8}\label{eq:lambda}\\
  \mu &= \sqrt{\frac{a_0}{K}}
    \bigl(2\,\lambda_1 + 2\,\lambda_2\sin\theta + 2\,\lambda_3\cos\theta\bigr),
    \tag{9}\label{eq:mu}\\
  \nu &= \sqrt{\frac{a_0}{K}}
    \bigl(\lambda_4\sin\theta + \lambda_5\cos\theta\bigr).
    \tag{10}\label{eq:nu}
\end{align}
The factor $\sqrt{a_0/K}$ arises from the time-to-fuel scaling in the
state equations.  The maximum principle requires that the thrust
direction maximise $H$ at every impulse.  Therefore the thrust points
in the direction of the primer vector, and the optimum impulses occur
at longitudes $\theta$ where $p = \sqrt{\lambda^2 + \mu^2 + \nu^2}$
attains a (global) maximum.

\paragraph{Elliptic form of the primer locus.}  The components
\eqref{eq:lambda}--\eqref{eq:nu} depend on $\theta$ through sine and
cosine.  Following Lawden we combine $\lambda_2$ and $\lambda_3$ into
a magnitude and phase via
\[
  \sqrt{\lambda_2^2 + \lambda_3^2} =: L,\qquad
  \tan \theta_0 = \frac{\lambda_2}{\lambda_3}.
\]
Using the trigonometric identities
\begin{align}
  \lambda_2\cos\theta - \lambda_3\sin\theta &= L\sin(\theta - \theta_0),\\
  \lambda_2\sin\theta + \lambda_3\cos\theta &= L\cos(\theta - \theta_0),
\end{align}
we may rewrite Eqs.~(\ref{eq:lambda})--(\ref{eq:nu}) as
\begin{align}
  \lambda(\theta) &= \sqrt{\frac{a_0}{K}}\,L\,\sin(\theta - \theta_0),
    \tag{11}\label{eq:lambdatheta}\\
  \mu(\theta) &= \sqrt{\frac{a_0}{K}}
    \bigl(2\,\lambda_1 + 2 L\cos(\theta - \theta_0)\bigr),
    \tag{12}\label{eq:mutheta}\\
  \nu(\theta) &= \sqrt{\frac{a_0}{K}}
    \Bigl( \frac{\lambda_3\lambda_4 - \lambda_2\lambda_5}{L}
           \sin(\theta - \theta_0)
           + \frac{\lambda_2\lambda_4 + \lambda_3\lambda_5}{L}
           \cos(\theta - \theta_0)\Bigr).
    \tag{13}\label{eq:nutheta}
\end{align}
In this form $(\lambda,\mu,\nu)$ traces an ellipse in the three-
dimensional space spanned by the radial, tangential and normal axes as
the true longitude $\theta$ varies.  Two special cases of this
ellipse lead to multiple maxima of the primer magnitude $p$ and
therefore to multi-impulse solutions.

\section*{3\quad Classification of primer loci}

Multiple-impulse solutions require that the primer magnitude $p$ have
two or more equal maxima as a function of $\theta$.  The geometry of
\eqref{eq:lambdatheta}--\eqref{eq:nutheta} shows that this can occur
only when the ellipse of the primer locus satisfies one of three
conditions.  We summarise these cases before solving the associated
two-point boundary-value problems.

\subsection*{3.1\quad Nodal case}

When the centre of the ellipse passes through the origin of the
$(\lambda,\mu,\nu)$-space, the primer magnitude is symmetric about
two diametrically opposite points on the major axis of the ellipse.
These points differ by $\pi$ in longitude.  Algebraically, the
condition that the ellipse be centred at the origin is that
$\lambda_1 = 0$.  With that assumption, the maxima of
\eqref{eq:lambdatheta}--\eqref{eq:nutheta} occur at $\theta_1$ and
$\theta_2 = \theta_1 + \pi$.  Equations
\eqref{eq:lambdatheta}--\eqref{eq:nutheta} then imply
\[
  \lambda(\theta_2) = -\lambda(\theta_1),\quad
  \mu(\theta_2)    = -\mu(\theta_1),\quad
  \nu(\theta_2)    = -\nu(\theta_1).
\]
Hence the two optimal impulses have equal magnitudes and opposite
components.

\subsection*{3.2\quad Non-degenerate case}

In the second case the ellipse passes through the $\mu$-axis (i.e.
the tangential axis).  This occurs when $\lambda_2\lambda_4 =
-\lambda_3\lambda_5$.  The primer locus then lies in a single plane
and has two symmetric maxima located at longitudes $\theta_1$ and
$\theta_2$ equidistant from the phase angle $\theta_0$, i.e. $\theta_2
- \theta_0 = \theta_0 - \theta_1$.  At these points one finds
\[
  \lambda(\theta_2) = -\lambda(\theta_1),\quad
  \mu(\theta_2)    =  \mu(\theta_1),\quad
  \nu(\theta_2)    = -\nu(\theta_1).
\]
The algebraic conditions defining the non-degenerate case will be
revisited in Section~\ref{sec:nondegenerate-bvp}.

\subsection*{3.3\quad Singular case}

When the ellipse degenerates into a circle, the primer magnitude is
constant for all $\theta$.  This singular case arises when
$\lambda_1 = 0$, $\lambda_2\lambda_4 = -\lambda_3\lambda_5$ and
$3 (\lambda_2^2 + \lambda_3^2)^2 = (\lambda_3\lambda_4 -
\lambda_2\lambda_5)^2$.  Up to scaling one finds
\[
  \lambda(\theta) = \tfrac12\sin(\theta - \theta_0),\quad
  \mu(\theta)    =      \cos(\theta - \theta_0),\quad
  \nu(\theta)    = \tfrac{\sqrt{3}}{2}\sin(\theta - \theta_0).
\]
Because every point on the orbit yields the same primer magnitude, the
impulses can be placed anywhere; the only requirement is that the
net change in the orbital elements satisfy the boundary conditions.
Section~\ref{sec:singular-case} presents the associated solution.

\section*{4\quad Two-point boundary-value problem: nodal case}

In the nodal case $\lambda_1 = 0$, so the centre of the primer ellipse
is at the origin and the two impulse locations are separated by
$180^\circ$ in longitude.  We choose the $x$-axis along the line of nodes of
the transfer, so that the impulses occur at $\theta_1 = 0$ and
$\theta_2 = \pi$.  Let $u_1$ and $u_2$ denote the magnitudes of the
two impulses and define $u = u_1 + u_2$ and $\delta = u_2/u$.  The
state equations~\eqref{eq:dx1du}--\eqref{eq:dx5du} are linear in the
control and may be integrated piecewise over the two impulses.  Using
the symmetry of the primer vector one finds for the net changes in the
orbital elements
\begin{align}
  \Delta x_1 &= 2\,u\,\sqrt{\frac{a_0}{K}}\,l_{T,1}\,(1-2\delta),
    \tag{A1}\label{eq:nodal-dx1}\\
  \Delta x_2 &= -u\,\sqrt{\frac{a_0}{K}}\,l_{R,1},
    \tag{A2}\label{eq:nodal-dx2}\\
  \Delta x_3 &= 2\,u\,\sqrt{\frac{a_0}{K}}\,l_{T,1},
    \tag{A3}\label{eq:nodal-dx3}\\
  \Delta x_4 &= 0,
    \tag{A4}\label{eq:nodal-dx4}\\
  \Delta x_5 &= u\,\sqrt{\frac{a_0}{K}}\,l_{N,1}.
    \tag{A5}\label{eq:nodal-dx5}
\end{align}
Here $(l_{R,1},l_{T,1},l_{N,1})$ denotes the thrust direction of the
first impulse; the second impulse has direction $-(l_{R,1},l_{T,1},l_{N,1})$ because the primer vector reverses when $\theta$ is increased by~$\pi$.  Equations~\eqref{eq:nodal-dx1}--\eqref{eq:nodal-dx5} follow directly from integrating Eqs.~\eqref{eq:dx1du}--\eqref{eq:dx5du} over the two impulses and making use of the fact that $\sin(\theta_1)=0$, $\cos(\theta_1)=1$ and $\sin(\theta_2)=0$, $\cos(\theta_2)=-1$.

Solving Eqs.~(\ref{eq:nodal-dx1})--\eqref{eq:nodal-dx5} for the total
impulse $u$ and the split $u_1$ yields
\begin{align}
  u &= \sqrt{\frac{K}{a_0}}\,\sqrt{\Delta i^2 + \frac{1}{4}\,\Delta e_x^2 + \Delta e_y^2},
    \tag{A6}\label{eq:nodal-u}\\
  u_1 &= \frac{\Delta e_x + \Delta x_1}{2\,\Delta e_x}\,u.
    \tag{A7}\label{eq:nodal-u1}
\end{align}
Equation~\eqref{eq:nodal-u} states that the minimum total impulse for
a nodal transfer depends on the final change in inclination and
eccentricity but not on the change in semi-major axis $\Delta x_1$
(because $\lambda_1=0$).  The impulse split \eqref{eq:nodal-u1}
reduces to $u_1 = u/2$ when there is no change in the semi-major axis.

\paragraph{Admissibility conditions.}  Two inequalities limit the
region of applicability of the nodal solution.  First, the primer
vector at the location of an impulse must point in the same direction
as the control, which implies that $\Delta x_1^2 \leq \Delta e_x^2$.
Secondly, the maxima at $\theta_1$ and $\theta_2$ must be genuine
maxima of $p$; this requirement leads to $\Delta i^2 \geq 3\,\Delta
e_y^2$.  Transfers violating either of these inequalities lie in the
non-degenerate or singular regions and must be treated separately.

\paragraph{Optimal impulse directions.}  The Lagrange multipliers can
be obtained by differentiating the payoff ($u$) with respect to the
terminal state.  When the primer is normalised so that $p=1$ at the
impulse, the adjoints are
\begin{align}
  \lambda_1 &= 0,
    \tag{A8}\label{eq:nodal-l1}\\
  \lambda_2 &= \sqrt{\frac{K}{a_0}}\,\frac{\Delta e_y}{\sqrt{\Delta i^2 + \tfrac14\Delta e_x^2 + \Delta e_y^2}},
    \tag{A9}\label{eq:nodal-l2}\\
  \lambda_3 &= \sqrt{\frac{K}{a_0}}\,\frac{\Delta e_x}{4\,\sqrt{\Delta i^2 + \tfrac14\Delta e_x^2 + \Delta e_y^2}},
    \tag{A10}\label{eq:nodal-l3}\\
  \lambda_4 &= -\frac{3}{4}\,\sqrt{\frac{K}{a_0}}\,\frac{\Delta e_x\,\Delta e_y}{\Delta i}\,\frac{1}{\sqrt{\Delta i^2 + \tfrac14\Delta e_x^2 + \Delta e_y^2}},
    \tag{A11}\label{eq:nodal-l4}\\
  \lambda_5 &= \sqrt{\frac{K}{a_0}}\,\frac{\Delta i}{\sqrt{\Delta i^2 + \tfrac14\Delta e_x^2 + \Delta e_y^2}}.
    \tag{A12}\label{eq:nodal-l5}
\end{align}
Note that in Eq.~\eqref{eq:nodal-l4} the denominator involves the
magnitude $\Delta i$ of the inclination change.  The scanned
article displays $\Delta i_x$ in that denominator, which is
misleading because in the nodal case $\Delta i_y = 0$ and hence
$\Delta i = \Delta i_x$.  We retain the correct dependence on
$\Delta i$ for clarity.

The primer vector at the impulse is obtained by substituting
\eqref{eq:nodal-l2}--\eqref{eq:nodal-l5} into
\eqref{eq:lambdatheta}--\eqref{eq:nutheta} and setting $\theta=0$.
Normalising to unit magnitude one finds the thrust direction cosines
\begin{equation}
  l_{R,1} = -\frac{\Delta e_y}{\sqrt{\Delta i^2 + \tfrac14\Delta e_x^2 + \Delta e_y^2}},\quad
  l_{T,1} = \frac{\Delta e_x}{2\,\sqrt{\Delta i^2 + \tfrac14\Delta e_x^2 + \Delta e_y^2}},\quad
  l_{N,1} = \frac{\Delta i}{\sqrt{\Delta i^2 + \tfrac14\Delta e_x^2 + \Delta e_y^2}}.
  \tag{A13}\label{eq:nodal-dircos}
\end{equation}
These direction cosines are the same for the second impulse except
that the radial and normal components change sign because the
longitudes are separated by $\pi$.

\section*{5\quad Two-point boundary-value problem: non-degenerate case}
\label{sec:nondegenerate-bvp}

In the non-degenerate case the primer ellipse is not centred at the
origin but still yields two equal maxima at symmetrically located
longitudes.  Solving the boundary-value problem requires more
algebra.  Following Edelbaum we introduce an alternative normalisation
of the adjoints in which the combination
\[
  \alpha := \sqrt{\frac{a_0}{K}}\,\frac{\lambda_3\lambda_4 - \lambda_2\lambda_5}{\sqrt{\lambda_2^2 + \lambda_3^2}}
\]
is set equal to one.  Define further
\[
  R := \sqrt{\frac{a_0}{K}}\,\sqrt{\lambda_2^2 + \lambda_3^2},\quad
  C := \cos(\theta - \theta_0),\quad S := \sin(\theta - \theta_0).
\]
With these definitions the primer components become rational
functions of $R$, $C$ and $S$.  Eliminating $\lambda_1$ using the
analogue of Eq.~(21) leads to the thrust direction cosines at the
impulse points:
\begin{align}
  l_R &= \frac{\lambda}{p}
        = \frac{2R^2\,S}{\sqrt{1+R^2}\,\sqrt{4R^2 + (1-3R^2)C^2}},
        \tag{B1}\label{eq:nondeg-lR}\\
  l_T &= \frac{\mu}{p}
        = \frac{\sqrt{1+R^2}\,C}{\sqrt{4R^2 + (1-3R^2)C^2}},
        \tag{B2}\label{eq:nondeg-lT}\\
  l_N &= \frac{\nu}{p}
        = \frac{2R\,S}{\sqrt{1+R^2}\,\sqrt{4R^2 + (1-3R^2)C^2}}.
        \tag{B3}\label{eq:nondeg-lN}
\end{align}
Because the second impulse occurs at $\theta_2$ such that $\theta_2
-\theta_0 = -(\theta_1-\theta_0)$, the values of $C$ and $S$ at the
two impulses are $\pm C$ and $\mp S$, respectively.  However,
Eq.~\eqref{eq:nondeg-lT} shows that the tangential direction cosine
$l_T$ has the same sign at both impulses whereas the radial and normal
components reverse, consistent with the conditions defining the
non-degenerate case.

\paragraph{Change in orbital elements.}  Substituting
\eqref{eq:nondeg-lR}--\eqref{eq:nondeg-lN} into the state
equations and integrating over the two impulses yields the total
changes in the orbital elements (in the temporary reference frame with
$x$-axis aligned with $\theta_0$):
\begin{align}
  \frac{\Delta x_1}{u} &=
    \frac{2C}{\sqrt{4R^2 + (1-3R^2)C^2}}\,\sqrt{1+R^2},
    \tag{B4}\label{eq:nondeg-dx1}\\
  \frac{\Delta x_2}{u} &=
    \frac{2\,(C^2 + R^2)}{\sqrt{1+R^2}\,\sqrt{4R^2 + (1-3R^2)C^2}},
    \tag{B5}\label{eq:nondeg-dx2}\\
  \frac{\Delta x_3}{u} &=
    \frac{2C\,S\,(1-2\delta)}{\sqrt{1+R^2}\,\sqrt{4R^2 + (1-3R^2)C^2}},
    \tag{B6}\label{eq:nondeg-dx3}\\
  \frac{\Delta x_4}{u} &=
    \frac{2R\,S^2}{\sqrt{1+R^2}\,\sqrt{4R^2 + (1-3R^2)C^2}},
    \tag{B7}\label{eq:nondeg-dx4}\\
  \frac{\Delta x_5}{u} &=
    \frac{2R\,C\,S\,(1-2\delta)}{\sqrt{1+R^2}\,\sqrt{4R^2 + (1-3R^2)C^2}}.
    \tag{B8}\label{eq:nondeg-dx5}
\end{align}
The eccentricity change vector $(\Delta x_2,\Delta x_3)$ makes an
angle $\omega$ with respect to the $x$-axis, while the inclination
change vector $(\Delta x_4,\Delta x_5)$ makes an angle $\Omega$.
Define
\[
  \tan \omega = \frac{\Delta x_2}{\Delta x_3},\quad
  \tan \Omega = \frac{\Delta x_4}{\Delta x_5}.
\]
Eliminating $\delta$ from Eqs.~\eqref{eq:nondeg-dx2}--\eqref{eq:nondeg-dx5}
shows that the difference between these angles is related to $R$, $C$ and
$S$ via
\begin{equation}
  T := \tan(\omega - \Omega) =
  \frac{(1-2\delta)^2 C^2 S - (C^2 + R^2)S}{(1-2\delta)(1+R^2)\,C}.
  \tag{B9}\label{eq:T-def}
\end{equation}
Solving this equation for the impulse split yields
\begin{equation}
  1 - 2\delta =
    \frac{T(1+R^2) - \sqrt{T^2 (1+R^2)^2 + 4 S^2 (C^2 + R^2)}}{2 S C}.
  \tag{B10}\label{eq:delta-solve}
\end{equation}
Substituting \eqref{eq:delta-solve} back into
\eqref{eq:nondeg-dx2}--\eqref{eq:nondeg-dx5} and eliminating $S$ in
favour of $C$ through $S = \sqrt{1-C^2}$ (with the appropriate sign)
allows the total magnitudes of the eccentricity and inclination
changes to be expressed in closed form as
\begin{align}
  \frac{\Delta e}{u} &:= \frac{1}{u}\,\sqrt{(\Delta x_2)^2 + (\Delta x_3)^2}
    = \sqrt{\frac{4(C^2 + R^2) + 2 T^2 (1+R^2) - 2 T\,\sqrt{T^2 (1+R^2)^2 + 4 S^2 (R^2 + C^2)}}{4 R^2 (1-3R^2) C^2}},
    \tag{B11}\label{eq:Delta-e}\\
  \frac{\Delta i}{u} &:= \frac{1}{u}\,\sqrt{(\Delta x_4)^2 + (\Delta x_5)^2}
    = R\,\sqrt{\frac{4 S^2 + 2 T^2 (1+R^2) - 2 T\,\sqrt{T^2 (1+R^2)^2 + 4 S^2 (R^2 + C^2)}}{4 R^2 (1-3R^2) C^2}}.
    \tag{B12}\label{eq:Delta-i}
\end{align}
These expressions simplify further when one introduces the change in
semi-major axis.  By eliminating $C$ and $S$ one obtains a relation
between $R$ and the transfer parameters.  To this end, multiply
\eqref{eq:Delta-e} by $R$ and divide it by
\eqref{eq:nondeg-dx1}, and divide \eqref{eq:Delta-i} by
\eqref{eq:nondeg-dx1}.  Subtracting the squares of the resulting
equations yields
\begin{equation}
  C^2 = \frac{R^2 (1 - R^2)\,(\Delta a/a_0)^2}{2 R^2 (\Delta a/a_0)^2 - (1+R^2)
    \bigl(R^2 \Delta e^2 - \Delta i^2\bigr)}.
  \tag{B13}\label{eq:C2}
\end{equation}
Finally one solves Eq.~\eqref{eq:C2} together with the definitions of
$\Delta e$ and $\Delta i$ for $R$.  After some algebra one arrives at
\begin{equation}
  R^2 = 1 + \frac{1}{2}
    \left[\frac{(\Delta a/a_0)^2 + \Delta i^2 - \Delta e^2}
    {\Delta i\,\Delta e\,\sin(\omega - \Omega)}\right]^2
    -
    \sqrt{
      \left[\frac{(\Delta a/a_0)^2 + \Delta i^2 - \Delta e^2}
      {\Delta i\,\Delta e\,\sin(\omega - \Omega)}\right]^2
      + \frac{1}{4}
      \left[\frac{(\Delta a/a_0)^2 + \Delta i^2 - \Delta e^2}
      {\Delta i\,\Delta e\,\sin(\omega - \Omega)}\right]^4}
  .
  \tag{B14}\label{eq:R2}
\end{equation}
In the scanned article Eq.~(66) contains typographical errors in the
placement of parentheses; Eq.~\eqref{eq:R2} above is the corrected
expression.  It shows that $R$ lies between $0$ and $1$ because the
square root selects the negative branch.

\paragraph{Total impulse.}  The Hamiltonian at the impulse is
proportional to the primer magnitude $p$, which is normalised to one.
The resulting payoff (total $u$) is obtained by substituting the
solutions of the boundary-value problem back into the cost function.
Eliminating $R$, $C$ and $S$ yields
\begin{equation}
  u = \sqrt{\frac{K}{2a_0}}
      \sqrt{\Delta i^2 + \Delta e^2 - \frac{(\Delta a/a_0)^2}{2}
      + \sqrt{\bigl(\Delta i^2 - \Delta e^2 + (\Delta a/a_0)^2\bigr)^2
              + 4\,\Delta i^2\,\Delta e^2\,\sin^2(\omega-\Omega)}}.
  \tag{B15}\label{eq:nondeg-u-general}
\end{equation}
Rotating the temporary $x$-axis back to the line of nodes and
expressing $\Delta e$ in terms of its Cartesian components
$\Delta e_x$ and $\Delta e_y$ yields the equivalent expression
\begin{equation}
  u = \sqrt{\frac{K}{2a_0}}
      \sqrt{\Delta i^2 + \Delta e_x^2 + \Delta e_y^2
            - \frac{(\Delta a/a_0)^2}{2}
            + \sqrt{\bigl(\Delta i^2 - \Delta e_x^2 - \Delta e_y^2
                        + (\Delta a/a_0)^2\bigr)^2
                    + 4\,\Delta i^2\,\Delta e_y^2}}.
  \tag{B16}\label{eq:nondeg-u}
\end{equation}
Equation~\eqref{eq:nondeg-u} corrects a sign error in the original
source: the term inside the outer square root that depends on
$\Delta a$ appears with the factor $-\tfrac12$, not $+\tfrac12$.

\paragraph{Adjoint variables.}  As in the nodal case, the Lagrange
multipliers may be obtained by differentiating the payoff
\eqref{eq:nondeg-u} with respect to the terminal state.  After
straightforward differentiation and simplification one finds
\begin{align}
  \lambda_1 &= \frac{K\,\Delta a}{2a_0^2\,u}
    \left[-\frac{1}{2}
    + \frac{\Delta i^2 - \Delta e_x^2 - \Delta e_y^2 + (\Delta a/a_0)^2}{
      \sqrt{\bigl(\Delta i^2 - \Delta e_x^2 - \Delta e_y^2 + (\Delta a/a_0)^2\bigr)^2
      + 4\,\Delta i^2\,\Delta e_y^2}}
    \right],
    \tag{B17}\label{eq:nondeg-l1}\\
  \lambda_2 &= \frac{K\,\Delta e_y}{2a_0\,u}
    \left[1
    - \frac{-\Delta i^2 - \Delta e_x^2 - \Delta e_y^2 + (\Delta a/a_0)^2}{
      \sqrt{\bigl(\Delta i^2 - \Delta e_x^2 - \Delta e_y^2 + (\Delta a/a_0)^2\bigr)^2
      + 4\,\Delta i^2\,\Delta e_y^2}}
    \right],
    \tag{B18}\label{eq:nondeg-l2}\\
  \lambda_3 &= \frac{K\,\Delta e_x}{2a_0\,u}
    \left[1
    - \frac{\Delta i^2 - \Delta e_x^2 - \Delta e_y^2 + (\Delta a/a_0)^2}{
      \sqrt{\bigl(\Delta i^2 - \Delta e_x^2 - \Delta e_y^2 + (\Delta a/a_0)^2\bigr)^2
      + 4\,\Delta i^2\,\Delta e_y^2}}
    \right],
    \tag{B19}\label{eq:nondeg-l3}\\
  \lambda_4 &= -\frac{K\,\Delta e_y}{2a_0\,u}
    \left[\frac{2\,\Delta e_x\,\Delta i}{
      \sqrt{\bigl(\Delta i^2 - \Delta e_x^2 - \Delta e_y^2 + (\Delta a/a_0)^2\bigr)^2
      + 4\,\Delta i^2\,\Delta e_y^2}}
    \right],
    \tag{B20}\label{eq:nondeg-l4}\\
  \lambda_5 &= \frac{K\,\Delta i}{2a_0\,u}
    \left[1
    + \frac{\Delta i^2 - \Delta e_x^2 + \Delta e_y^2 + (\Delta a/a_0)^2}{
      \sqrt{\bigl(\Delta i^2 - \Delta e_x^2 - \Delta e_y^2 + (\Delta a/a_0)^2\bigr)^2
      + 4\,\Delta i^2\,\Delta e_y^2}}
    \right].
    \tag{B21}\label{eq:nondeg-l5}
\end{align}
These expressions correct several misprints in the scanned article,
notably the sign of the prefactor in $\lambda_1$ and the placement of
the terms inside the radicals.

\section*{6\quad Singular case}
\label{sec:singular-case}

When the primer locus becomes a circle, the magnitude of the primer
vector is independent of $\theta$.  From the conditions
\[
  \lambda_1 = 0,\quad \lambda_2\lambda_4 = -\lambda_3\lambda_5,\quad
  3(\lambda_2^2 + \lambda_3^2)^2 = (\lambda_3\lambda_4 - \lambda_2\lambda_5)^2,
\]
one obtains up to scaling
\[
  \lambda(\theta) = \tfrac12\sin(\theta - \theta_0),\quad
  \mu(\theta)    =      \cos(\theta - \theta_0),\quad
  \nu(\theta)    = \tfrac{\sqrt{3}}{2}\sin(\theta - \theta_0).
\]
Every value of $\theta$ is a maximum of $p$ so the impulses may be
distributed arbitrarily.  A convenient choice is to place the two
impulses at longitudes separated by $\pi$ and with equal magnitudes
$u/2$; however the solution is not unique.

\paragraph{Change in semi-major axis.}  In the singular region the
semi-major axis is no longer determined by the eccentricity and
inclination changes: the payoff does not depend on $a$ because
$\lambda_1=0$.  Setting $\lambda_1=0$ in the general expression
\eqref{eq:nondeg-l1} and solving for $\Delta a$ yields
\begin{equation}
  (\Delta a)^2 = a_0^2\Bigl( \Delta e_x^2 + \Delta e_y^2 +
    2\,\sqrt{3}\,\Delta i\,\Delta e_y - \Delta i^2 \Bigr).
  \tag{C1}\label{eq:singular-a}
\end{equation}
If the required $\Delta a$ of a transfer satisfies
$|\Delta a|\leq a_0\,\sqrt{\Delta e_x^2 + \Delta e_y^2 + 2\sqrt{3}
\,\Delta i\,\Delta e_y - \Delta i^2}$, then the transfer lies in the
singular or nodal region.  The payoff for the singular case follows
from Eqs.~\eqref{eq:nondeg-u} or \eqref{eq:nondeg-u-general} by
substituting the bound~\eqref{eq:singular-a} and simplifying:
\begin{equation}
  u = \sqrt{\frac{K}{4a_0}}
        \sqrt{ \bigl( \sqrt{3}\,\Delta i + \Delta e_y\bigr)^2 + \Delta e_x^2 }.
  \tag{C2}\label{eq:singular-u}
\end{equation}
This expression depends only on the eccentricity and inclination
changes; the semi-major axis change is limited by
\eqref{eq:singular-a} and does not affect $u$.

\paragraph{Orientation of the singular locus.}  The orientation
$\theta_0$ of the primer circle relative to the line of nodes is
obtained by substituting the derivatives of the payoff with respect to
$\Delta e_x$ and $\Delta e_y$ into Eq.~(14).  A compact formula is
\begin{equation}
  \tan \theta_0 = \frac{\Delta e_y + \sqrt{3}\,\Delta i}{\Delta e_x}.
  \tag{C3}\label{eq:singular-theta0}
\end{equation}
This shows that when the eccentricity and inclination change vectors
are aligned ($\Delta e_x=0$), the primer circle is oriented midway
between them.

\paragraph{Reduced dynamics.}  In the singular region the effective
dimension of the reachable set is reduced.  Choosing the $x$-axis
along $\theta_0$ one may substitute the singular expressions for
$(\lambda,\mu,\nu)$ into the state equations.  The resulting rates of
change of the elements are
\begin{align}
  \frac{\mathrm{d}x_1}{\mathrm{d}u} &= 2C,
    \tag{C4}\label{eq:singular-dx1}\\
  \frac{\mathrm{d}x_2}{\mathrm{d}u} &= \frac{3}{2}\,S C,
    \tag{C5}\label{eq:singular-dx2}\\
  \frac{\mathrm{d}x_3}{\mathrm{d}u} &= 2 - \frac{3}{2}\,S^2,
    \tag{C6}\label{eq:singular-dx3}\\
  \frac{\mathrm{d}x_4}{\mathrm{d}u} &= \frac{\sqrt{3}}{2}\,S^2,
    \tag{C7}\label{eq:singular-dx4}\\
  \frac{\mathrm{d}x_5}{\mathrm{d}u} &= \frac{\sqrt{3}}{2}\,S C.
    \tag{C8}\label{eq:singular-dx5}
\end{align}
These relations show that $x_2$ is proportional to $x_5$ and $x_3$ is
proportional to $x_4$ along a single impulse.  By using two impulses
with different values of $S^2$ and appropriate splits one can reach
any point inside the convex hull defined by the envelope of these
curves; the details of these constructions lie beyond the scope of
this derivation.

\section*{7\quad Discussion and conclusions}

The preceding derivations provide complete analytic expressions for
minimum-impulse transfers between nearby orbits under the assumption
that both eccentricity and inclination remain small.  Three classes of
solutions arise: the nodal case, in which the two impulses occur
exactly on the line of nodes; the non-degenerate case, in which the
primer locus is an off-centre ellipse and the impulses occur away
from the nodes with unequal splits; and the singular case, in which
the primer locus is a circle and any distribution of thrust along the
optimal direction yields the same fuel consumption.  Correcting
typographical errors in the original source, we presented explicit
expressions for the total impulse, the impulse split, the thrust
direction cosines and the associated adjoint variables in each case.

The singular solutions are of particular interest because they allow
finite-thrust arcs to replace sequences of impulses without loss of
optimality.  They also hint at the existence of nonlinear singular
solutions in the full (nonlinear) optimal control problem.  Future
research may extend the present linear analysis to include higher
order effects, finite thrust arcs and terminal constraints on
position.

\end{document}
